\DocumentMetadata{
  pdfversion=2.0,
  pdfstandard=ua-2,
  lang=en-US,
  tagging=on,
  tagging-setup={math/setup=mathml-SE,tabsorder=structure}
}
\documentclass[11pt,letterpaper]{article}
\usepackage[margin=0.68in,includefoot]{geometry}
\usepackage{newcomputermodern}
\usepackage{amsmath,amscd,mathtools}
\usepackage{tikz,adjustbox}
\providecommand{\square}{\mdlgwhtsquare}
\usepackage[hidelinks]{hyperref}
\hypersetup{
  pdftitle={Numerical Analysis First-Year Exam, May 2023},
  pdfauthor={University of Florida Department of Mathematics},
  pdfsubject={Graduate mathematics examination},
  pdfdisplaydoctitle=true,
  bookmarksopen=true
}
\setcounter{secnumdepth}{0}
\setlength{\parindent}{0pt}
\setlength{\parskip}{0.28em}
\setlength{\emergencystretch}{3em}
\setlength{\leftmargini}{2.15em}
\setlength{\leftmarginii}{2.65em}
\setlength{\leftmarginiii}{3em}
\pagestyle{empty}
\raggedbottom
\allowdisplaybreaks
\NewTaggingSocketPlug{block/list/label}{examproblem}{%
  \tagstructbegin{tag=Lbl}\tagstructend%
  \tagstructbegin{tag=\LBody}%
  \tagstructbegin{tag=\ProblemHeadingTag,title-o={\ProblemHeadingText}}%
  \tagmcbegin{tag=\ProblemHeadingTag,actualtext={\ProblemHeadingText}}%
  #2%
  \tagmcend%
  \tagstructend%
}
\newenvironment{examproblems}[2]{%
  \def\ProblemHeadingTag{#2}%
  \AssignTaggingSocketPlug{block/list/label}{examproblem}%
  \begin{enumerate}%
  \setcounter{enumi}{#1}%
  \renewcommand{\labelenumi}{\textbf{\theenumi.}}%
  \setlength{\topsep}{0.35em}%
  \setlength{\partopsep}{0pt}%
  \setlength{\itemsep}{0.48em}%
  \setlength{\parsep}{0.18em}%
}{\end{enumerate}}
\newenvironment{examsubparts}{%
  \AssignTaggingSocketPlug{block/list/label}{default}%
  \begin{enumerate}%
  \setlength{\topsep}{0.18em}%
  \setlength{\partopsep}{0pt}%
  \setlength{\itemsep}{0.16em}%
  \setlength{\parsep}{0.1em}%
}{\end{enumerate}}
\newenvironment{examinstructions}{%
  \par\begingroup\itshape
}{%
  \par\endgroup\vspace{0.18em}
}
\makeatletter
\renewcommand\subsection{\@startsection{subsection}{2}{0pt}%
  {-1.25ex plus -.25ex minus -.1ex}%
  {0.55ex plus .1ex}%
  {\normalfont\large\bfseries}}
\renewcommand{\@maketitle}{%
  \newpage\null\vskip -1em%
  {\footnotesize\bfseries
    \href{https://www.ufl.edu/}{University of Florida}, \href{https://math.ufl.edu/}{Department of Mathematics}\par}%
  \vskip -0.5em%
  \tagstructbegin{tag=H1,title={Numerical Analysis First-Year Exam, May 2023}}%
  \tagpdfparaOff%
  \tagmcbegin{tag=H1}%
    {\LARGE\bfseries\href{https://gma.math.ufl.edu/exams/numerical-analysis-first-year/}{Numerical Analysis First-Year Exam}, May 2023\par}%
  \tagmcend%
  \tagpdfparaOn%
  \tagstructend%
  \vskip 0.25em%
  \tagmcbegin{artifact}%
    \hrule height 1pt%
  \tagmcend%
  \vskip 1em%
}
\makeatother
\title{Numerical Analysis First-Year Exam, May 2023}
\author{}
\date{}
\begin{document}
\maketitle
\thispagestyle{empty}
\begin{examinstructions}
Do 4 (four) problems.
\end{examinstructions}
\begin{examproblems}{0}{H2}
\def\ProblemHeadingText{Problem 1.}
\item
Consider
\[
y'=4ty,\qquad t\in[0,1],\qquad y(0)=2, \tag{1}
\]
 which has solution \(Y(t)=2e^{2t^2}\).
\begin{examsubparts}
\item[\textnormal{(a)}]
Derive an error bound for the Euler scheme.
\item[\textnormal{(b)}]
Derive the Taylor method of order 2 for (1).
\end{examsubparts}
\def\ProblemHeadingText{Problem 2.}
\item
Find \(a,b,c\) so that the quadrature formula
\[
\int_0^2 f(x)\,dx=af(0)+bf(1)+cf(2)
\]
 has the largest degree of precision.
\def\ProblemHeadingText{Problem 3.}
\item
Suppose \(f\in C^{n+1}[a,b]\), and let \(p\in\mathcal{P}_n\) be a polynomial that interpolates the data \(\{(x_i,f(x_i))\}_{i=0}^n\), where \(x_0,\ldots,x_n\) are distinct points in \([a,b]\). Consider an arbitrary fixed \(x\in[a,b]\), and derive an exact expression for the error \(f(x)-p(x)\).
\def\ProblemHeadingText{Problem 4.}
\item
Let \(\mathcal{P}_1\) be the space of polynomials of degree at most one. Using the norm
\[
\lVert u\rVert_2=\left(\int_a^b u^2\,dx\right)^{1/2}.
\]
\begin{examsubparts}
\item[\textnormal{(a)}]
Find the least-squares approximation to \(f(x)=x^3\) in \(\mathcal{P}_1\) over \([a,b]=[-1,1]\).
\item[\textnormal{(b)}]
Find the least-squares approximation to \(f(x)=x^3\) in \(\mathcal{P}_1\) over \([a,b]=[0,1]\).
\end{examsubparts}
\def\ProblemHeadingText{Problem 5.}
\item
Let \(G=[0,2]\) and
\[
g(x)=\frac{1}{3}\left(\frac{x^3}{3}-x^2-\frac{5x}{4}+4\right).
\]
\begin{examsubparts}
\item[\textnormal{(a)}]
Use the contraction mapping theorem to prove that if \(x_0\in G\), then the sequence defined by \(x_{k+1}=g(x_k)\) \((k=0,1,\ldots)\) converges to a unique fixed point \(z\in G\).
\item[\textnormal{(b)}]
Consider the fixed point iteration method \(x_{k+1}=g(x_k)\) \((k=0,1,\ldots)\) for solving the nonlinear equation \(f(x)=0\). Consider choosing an iteration function of the form
\[
g(x)=x-af(x)-b(f(x))^2-c(f(x))^3,
\]
 where \(a,b,c\) are parameters to be determined. Assume~\(f\) is sufficiently differentiable and the corresponding iterations \(x_{k+1}=g(x_k)\) converge to a unique fixed point~\(z\). Find expressions for the parameters \(a,b,c\) in terms of functions of~\(z\) such that the iteration method is of fourth order.
\end{examsubparts}
\end{examproblems}
\end{document}
